\section{Quadratic logarithmic range on a periodic physical chain}
\label{sec:hastings_periodic_coarsening}

The grouping into supersites in
\cite[lines 609--639]{Hastings2006Solving} gives pair-operator bounds
proportional to the projection radius, rather than to the total number
of physical sites. This yields the quadratic logarithmic range in one
dimension. The construction below retains the physical periodic boundary
and includes arbitrary physical chain lengths. The hierarchy uses the
corrected deficiency and the additional closing window of
\cite{gap:hastings05_hierarchical_projection,gap:hastings05_hierarchical_energy}.

\begin{definition}[Regular periodic grouping]
    \label{def:hastings_regular_periodic_grouping}
    \lean{TNLean.MPS.AreaLaw.LocalSolution.periodicCoarseEdge,
        TNLean.MPS.AreaLaw.LocalSolution.periodicPhysicalCoordinate,
        TNLean.MPS.AreaLaw.LocalSolution.periodicPhysicalDistance,
        TNLean.MPS.AreaLaw.LocalSolution.periodicCoarseOperator,
        TNLean.MPS.AreaLaw.LocalSolution.periodicCoarseProjection}
    \leanok
    Label the sites by $(k,s)$ with $0\le k<N$ and $0\le s<2\ell$.
    Their physical coordinate is $2\ell k+s$, with cyclic distance on
    $2\ell N$ sites. Assign $(k,s)$ to edge $k$ when $s<\ell$, and to
    edge $k+1$ modulo $N$ otherwise. The operator $A_j$ is the sum of
    terms assigned to edge $j$. The periodic hierarchy is formed from
    these operators, with centred norm bound $K=4\ell J$.
\end{definition}

\begin{theorem}[Sharp regular grouping]
    \label{thm:hastings_regular_periodic_grouping}
    \lean{TNLean.MPS.AreaLaw.LocalSolution.card_periodicCoarseEdge_fibre,
        TNLean.MPS.AreaLaw.LocalSolution.finRotate_val_eq_succ_or_zero,
        TNLean.MPS.AreaLaw.LocalSolution.periodicPhysicalCoordinate_lt,
        TNLean.MPS.AreaLaw.LocalSolution.periodicCoarseEdge_internal_centre,
        TNLean.MPS.AreaLaw.LocalSolution.periodicCoarseEdge_support,
        TNLean.MPS.AreaLaw.LocalSolution.periodicCoarseEdge_zero_support,
        TNLean.MPS.AreaLaw.LocalSolution.periodicCoarseEdge_support_all,
        TNLean.MPS.AreaLaw.LocalSolution.sum_periodicCoarseOperator,
        TNLean.MPS.AreaLaw.LocalSolution.periodicCoarseOperator_mem_supportedOperators,
        TNLean.MPS.AreaLaw.LocalSolution.isHermitian_periodicCoarseOperator}
    \leanok
    \uses{def:hastings_regular_periodic_grouping}
    Every edge receives exactly $2\ell$ terms. A cyclic ball of radius
    $\ell$ about a centre assigned to $j$ lies in blocks $j-1,j$,
    including the last and first blocks when $j=0$. Thus cyclic-local
    Hermitian terms give Hermitian adjacent-pair operators, and
    $\sum_jA_j=\sum_iM_i$. The cyclic successor increases a natural
    block label by one, except that the last label returns to zero.
\end{theorem}
\begin{proof}\leanok
    For a fixed within-block coordinate $s$, exactly one block gives
    assignment to $j$: it is $j$ in the first half, and $j-1$ in the
    second half. This is a bijection between the centres assigned to
    $j$ and the $2\ell$ possible values of $s$.
    At an internal edge, the physical centre $z$ satisfies
    $2\ell j\le z+\ell$ and $z<2\ell j+\ell$. Its cyclic locality ball
    therefore lies between $2\ell(j-1)$ and $2\ell(j+1)$.
    The closing edge has centres in the first or last half-block, so
    the same assertion holds across the periodic cut. Summing the
    assigned terms proves the sum and Hermiticity assertions.
\end{proof}

\begin{definition}[Grouping an arbitrary periodic chain]
    \label{def:hastings_rounded_periodic_grouping}
    \lean{TNLean.MPS.AreaLaw.LocalSolution.roundedPeriodicBlock,
        TNLean.MPS.AreaLaw.LocalSolution.roundedPeriodicEdge,
        TNLean.MPS.AreaLaw.LocalSolution.roundedPeriodicDistance,
        TNLean.MPS.AreaLaw.LocalSolution.roundedPeriodicBlockCount,
        TNLean.MPS.AreaLaw.LocalSolution.roundedPeriodicOperator,
        TNLean.MPS.AreaLaw.LocalSolution.roundedPeriodicProjection}
    \leanok
    For $M$ physical sites and an integer $\ell\ge1$, put
    $N=\max\{1,\lfloor M/(2\ell)\rfloor\}$. The block containing
    $z\in\{0,\ldots,M-1\}$ has label
    $b(z)=\min\{\lfloor z/(2\ell)\rfloor,N-1\}$. Assign $z$ to
    edge $\lfloor(z+\ell)/(2\ell)\rfloor$ if this is less than $N$,
    and to edge zero otherwise. Distance remains the physical cyclic
    distance $\min\{|x-z|,M-|x-z|\}$. The pair operator $A_j$ sums the
    assigned terms, and its hierarchy uses $K=8\ell J$.
\end{definition}

\begin{theorem}[Rounded cyclic supports and the number of terms]
    \label{thm:hastings_rounded_periodic_grouping}
    \lean{TNLean.MPS.AreaLaw.LocalSolution.roundedPeriodicBlock_lt,
        TNLean.MPS.AreaLaw.LocalSolution.roundedPeriodicEdge_internal_centre,
        TNLean.MPS.AreaLaw.LocalSolution.roundedPeriodicEdge_support,
        TNLean.MPS.AreaLaw.LocalSolution.roundedPeriodicBlockCount_pos,
        TNLean.MPS.AreaLaw.LocalSolution.roundedPeriodicBlockCount_cover,
        TNLean.MPS.AreaLaw.LocalSolution.roundedPeriodicBlockCount_upper,
        TNLean.MPS.AreaLaw.LocalSolution.roundedPeriodicBlockCount_le,
        TNLean.MPS.AreaLaw.LocalSolution.roundedPeriodicEdge_support_of_cover,
        TNLean.MPS.AreaLaw.LocalSolution.roundedPeriodicEdge_zero_support,
        TNLean.MPS.AreaLaw.LocalSolution.card_roundedPeriodicEdge_fibre_le,
        TNLean.MPS.AreaLaw.LocalSolution.sum_roundedPeriodicOperator,
        TNLean.MPS.AreaLaw.LocalSolution.roundedPeriodicOperator_mem_supportedOperators,
        TNLean.MPS.AreaLaw.LocalSolution.isHermitian_roundedPeriodicOperator}
    \leanok
    \uses{def:hastings_rounded_periodic_grouping}
    The block count is positive, and is at most $M$ when $M\ge1$.
    Either $2\ell N\le M<2\ell(N+1)$ or the chain consists of one
    block. Every radius-$\ell$ cyclic ball lies in the pair assigned
    to its centre; the closing pair is the first and last blocks.
    Each edge receives at most $4\ell$ terms. Consequently cyclic-local
    Hermitian terms give Hermitian pair operators supported on adjacent
    blocks, with $\sum_jA_j=\sum_iM_i$.
\end{theorem}
\begin{proof}\leanok
    Division with remainder gives the block count and the bounds on
    $M$. The first $N-1$ blocks have width $2\ell$; the last contains
    the remainder and has width less than $4\ell$. An internal assigned
    centre again satisfies $2\ell j\le z+\ell$ and
    $z<2\ell j+\ell$, giving the adjacent pair. The closing assignment
    contains only centres in the first half-block or beyond the last
    internal boundary plus $\ell$. Cyclic locality then gives the first
    or last block. If $N=1$, every site already belongs to that block.
    For the number of centres, an internal assignment lies in an interval
    of length $2\ell$. For edge zero, the first $\ell$ centres and the
    remaining centres of the last block together number at most $4\ell$.
\end{proof}

\begin{theorem}[Local estimates give linear pair estimates]
    \label{thm:hastings_periodic_linear_pair_bounds}
    \lean{TNLean.MPS.AreaLaw.LocalSolution.norm_edgePartition_le_of_fibre_card,
        TNLean.MPS.AreaLaw.LocalSolution.residual_periodicCoarseOperator_le,
        TNLean.MPS.AreaLaw.LocalSolution.centeredNorm_periodicCoarseOperator_le,
        TNLean.MPS.AreaLaw.LocalSolution.residual_roundedPeriodicOperator_le,
        TNLean.MPS.AreaLaw.LocalSolution.centeredNorm_roundedPeriodicOperator_le}
    \leanok
    \uses{thm:hastings_regular_periodic_grouping,
        thm:hastings_rounded_periodic_grouping}
    Let the $M_i$ be Hermitian, let $\psi$ be unit, let $J,\varepsilon\ge0$, and suppose
    $\|M_i\|\le J$ and
    $\|(M_i-\langle\psi,M_i\psi\rangle)\psi\|\le\varepsilon$.
    In the rounded grouping,
    \begin{align*}
        \|A_j-\langle\psi,A_j\psi\rangle\|&\le8\ell J,\\
        \|(A_j-\langle\psi,A_j\psi\rangle)\psi\|&\le4\ell\varepsilon.
    \end{align*}
    In the regular grouping the respective constants are $4\ell J$
    and $2\ell\varepsilon$. More generally a sum over at most $B$ terms
    of norm at most $J$ has norm at most $BJ$.
\end{theorem}
\begin{proof}\leanok
    Sum the local residuals over the assigned centres. For the centred
    norm, use $|\langle\psi,A_j\psi\rangle|\le\|A_j\|$ and
    $\|A_j\|\le\sum_{i\mapsto j}\|M_i\|$, followed by the proved
    number of assigned centres.
\end{proof}

\begin{theorem}[Absorbing a linear truncation prefactor]
    \label{thm:hastings_periodic_linear_truncation_error}
    \lean{TNLean.MPS.AreaLaw.LocalSolution.hierarchyTruncationError_mul_length_le_eighth_gap,
        TNLean.MPS.AreaLaw.LocalSolution.card_le_two_mul_re_trace_of_periodicHierarchyProjection_of_error}
    \leanok
    \uses{thm:hastings05_explicit_projection_length,
        thm:hastings05_physical_gap_count}
    Let $h\ge2$, $J_*,\Delta>0$, and
    $r\ge512h^2(1+\max\{0,\log(J_*/\Delta)\})$. Then
    $2^hJ_*r e^{-r}\le\Delta/8$.
    In the quantitative periodic hierarchy, a centred pair bound
    $K\le J_*r$ and any physical approximation error
    $\delta\le\Delta/8$ therefore suffice for the count of at most
    twice the physical ground-space dimension.
\end{theorem}
\begin{proof}\leanok
    \uses{thm:hastings05_explicit_projection_length,
        thm:hastings05_physical_gap_count}
    Apply the window-energy bound with $K=2J_*r/9$, and use
    $e^{-r}\le e^{-r/(4h)}$. The sum of all radii is at most
    $\Delta/4$; adding the two expectation errors costs at most
    $2\delta\le\Delta/4$. The physical spectral gap then gives the count.
\end{proof}

\begin{theorem}[Explicit integer projection radius]
    \label{thm:hastings_periodic_chosen_radius}
    \lean{TNLean.MPS.AreaLaw.LocalSolution.hierarchyProjectionLength_roundedPeriodic_le,
        TNLean.MPS.AreaLaw.LocalSolution.roundedPeriodicProjectionRadius,
        TNLean.MPS.AreaLaw.LocalSolution.roundedPeriodicProjectionRadius_pos,
        TNLean.MPS.AreaLaw.LocalSolution.hierarchyProjectionLength_le_chosen_roundedPeriodicRadius,
        TNLean.MPS.AreaLaw.LocalSolution.roundedPeriodicProjectionRadius_lt}
    \leanok
    Let $M\ge1$, $\eta>0$, and write
    \begin{align*}
        A_M&=512(\lceil\log_2 M\rceil+2)^2
            \bigl(1+\max\{0,\log((8J/\eta)/\Delta)\}\bigr),\\
        \ell&=\lceil A_M/\eta\rceil,\qquad r=\eta\ell.
    \end{align*}
    Then $\ell\ge1$ and $\ell<A_M/\eta+1$.
    For $N=\max\{1,\lfloor M/(2\ell)\rfloor\}$ and
    $h=\max\{2,\lceil\log_2N\rceil\}$, the actual hierarchy length
    $512h^2(1+\max\{0,\log((8J/\eta)/\Delta)\})$ is at most $r$.
\end{theorem}
\begin{proof}\leanok
    Since $N\le M$, the hierarchy depth is at most
    $\lceil\log_2M\rceil+2$. Thus its length is at most $A_M$.
    The inequalities $A_M/\eta\le\lceil A_M/\eta\rceil<A_M/\eta+1$
    give the assertions, including the length inequality for the block
    count determined by this same choice of $\ell$.
\end{proof}

\begin{theorem}[The hierarchy from periodic Gaussian-localization estimates]
    \label{thm:hastings_periodic_physical_log_squared}
    \lean{TNLean.MPS.AreaLaw.LocalSolution.exists_unit_mem_roundedPeriodicProjection,
        TNLean.MPS.AreaLaw.LocalSolution.card_le_two_mul_re_trace_of_roundedPeriodicProjection,
        TNLean.MPS.AreaLaw.LocalSolution.exists_unit_mem_roundedPeriodicProjection_of_chosen_radius,
        TNLean.MPS.AreaLaw.LocalSolution.card_le_two_mul_re_trace_of_roundedPeriodicProjection_of_chosen_radius}
    \leanok
    \uses{thm:hastings_periodic_chosen_radius,
        thm:hastings_periodic_linear_pair_bounds,
        thm:hastings05_quantitative_periodic}
    Let $H=\sum_iH_i$ on a periodic chain of $M\ge1$ sites, with
    Hermitian $H_i$, spectral gap $\Delta>0$, and a unit ground vector
    $\psi$ of energy $E_0$. Let $J,\eta>0$ and choose $\ell,r$ as in
    Theorem~\ref{thm:hastings_periodic_chosen_radius}.
    Suppose Hermitian $M_i$ act in the cyclic radius-$\ell$ balls and satisfy
    \begin{align*}
        \|M_i\|&\le J,\\
        \|(M_i-\langle\psi,M_i\psi\rangle)\psi\|&\le Je^{-r},\\
        \|\widetilde H_i-M_i\|&\le Je^{-r},
    \end{align*}
    where $\widetilde H_i$ is the full Gaussian filter of $H_i$ under $H$.
    Then the corrected periodic hierarchy constructed from the $M_i$
    retains a unit vector $\phi$ at every tree node, with
    $\|\phi-\psi\|\le2(3^{h+1}-1)e^{-r/8}$.
    Every finite orthonormal family in its final retained space has at
    most $2\dim\ker(H-E_0)$ members. The residual hypothesis is equation
    (low); the third hypothesis is the preceding Gaussian-localization
    estimate, rather than an estimate comparing $H_i$ directly with $M_i$.
    More generally, the same conclusions hold at any integer $\ell\ge1$
    with $r=\eta\ell$ satisfying the explicit length inequality for its
    own hierarchy depth.
\end{theorem}
\begin{proof}\leanok
    \uses{thm:hastings_periodic_linear_pair_bounds,
        thm:hastings_periodic_linear_truncation_error,
        thm:hastings_periodic_chosen_radius,
        thm:hastings05_quantitative_periodic}
    The derived pair bound is $K=8\ell J=(8J/\eta)r$.
    The local residuals give a pair residual at most $K e^{-r}$, so the
    quantitative hierarchy produces the nearby retained unit vector.
    The full filtered terms sum exactly to $H$. Hence
    $\|H-\sum_iM_i\|\le MJe^{-r}$, and $M\le4\ell N$ bounds this by
    $2^h(8J/\eta)r e^{-r}\le\Delta/8$. Apply the physical-gap count.
\end{proof}

\begin{theorem}[The sharper constant for regular periodic blocks]
    \label{thm:hastings_regular_periodic_log_squared}
    \lean{TNLean.MPS.AreaLaw.LocalSolution.exists_unit_mem_periodicCoarseProjection,
        TNLean.MPS.AreaLaw.LocalSolution.card_le_two_mul_re_trace_of_periodicCoarseProjection,
        TNLean.MPS.AreaLaw.LocalSolution.hierarchyProjectionLength_periodicCoarse_le}
    \leanok
    For $M=2\ell N$ sites in regular blocks, the same retained-vector
    and counting assertions hold with $K=4\ell J$ and effective coupling
    $4J/\eta$. The length threshold is at most
    $512(\lceil\log_2M\rceil+2)^2
    (1+\max\{0,\log((4J/\eta)/\Delta)\})$.
    For the count, any comparison terms $G_i$ with sum $H$ and
    $\|G_i-M_i\|\le Je^{-r}$ suffice; in the physical application these
    are the full filtered terms.
\end{theorem}
\begin{proof}\leanok
    \uses{thm:hastings_regular_periodic_grouping,
        thm:hastings_periodic_linear_pair_bounds,
        thm:hastings05_quantitative_periodic}
    Use the exact number $2\ell$ of terms assigned to each edge in place
    of $4\ell$. The remaining argument is unchanged.
\end{proof}

\begin{theorem}[The metric of the original physical cycle]
    \label{thm:hastings_periodic_physical_metric}
    \lean{TNLean.MPS.AreaLaw.LocalSolution.roundedPeriodicDistance_eq_valMinAbs,
        TNLean.MPS.AreaLaw.LocalSolution.roundedPeriodicDistance_triangle,
        TNLean.MPS.AreaLaw.LocalSolution.periodicPhysicalMetricSpace,
        TNLean.MPS.AreaLaw.LocalSolution.card_periodicPhysical_metricRangeBall_le}
    \leanok
    On a cycle of $M\ge1$ sites, the distance
    $d(i,j)=\min\{|i-j|,M-|i-j|\}$ is the least absolute residue of
    $i-j$ modulo $M$. It is a metric, and every closed radius-$R$ ball
    has at most $2\lceil R\rceil+1$ sites.
\end{theorem}
\begin{proof}\leanok
    The least absolute residue of a sum has absolute value no larger
    than the sum of the absolute values of the two representatives.
    This proves the triangle inequality; symmetry and separation are
    immediate. Subtract the centre of a ball and take least absolute
    representatives. This is an injection into the integer interval
    $[-\lceil R\rceil,\lceil R\rceil]$, which has $2\lceil R\rceil+1$
    elements.
\end{proof}

\begin{theorem}[The hierarchy from the original periodic Hamiltonian]
    \label{thm:hastings_periodic_physical_hierarchy}
    \lean{TNLean.MPS.AreaLaw.LocalSolution.physicalPeriodicProjectionRadius,
        TNLean.MPS.AreaLaw.LocalSolution.physicalPeriodicProjectionRadius_spec,
        TNLean.MPS.AreaLaw.LocalSolution.hierarchyProjectionLength_le_physicalPeriodicProjectionRadius,
        TNLean.MPS.AreaLaw.LocalSolution.physicalPeriodicProjectionRadius_lt,
        TNLean.MPS.AreaLaw.LocalSolution.periodicPhysicalTruncatedTerm,
        TNLean.MPS.AreaLaw.LocalSolution.exists_periodicPhysicalTruncatedTerm_localization,
        TNLean.MPS.AreaLaw.LocalSolution.exists_periodicPhysicalHierarchy}
    \leanok
    \uses{thm:hastings_periodic_physical_metric,
        thm:hastings_periodic_physical_log_squared,
        thm:hastings05_metric_facld,
        thm:hastings_zero_range_exact}
    Fix $R\ge0$, $J>0$ and $\Delta>0$. There are constants $c_1,\eta>0$
    and $C\ge0$, independent of physical chain length, with the following
    property. Let $H=\sum_iH_i$ be a Hermitian local Hamiltonian on a
    periodic chain of $M\ge1$ sites, with $\|H_i\|\le J$ and each $H_i$
    supported in the physical radius-$R$ ball about $i$. Let $P_0$ be
    its orthogonal ground projection, $E_0$ its ground energy, and
    $H-E_0\ge\Delta(1-P_0)$.
    Put $J_b=\max\{J,JC\}$ and
    \begin{align*}
        A_M&=512(\lceil\log_2M\rceil+2)^2
          \bigl(1+\max\{0,\log((8J_b/\eta)/\Delta)\}\bigr),\\
        \ell&=\max\{\lceil R\rceil,\lceil A_M/\eta\rceil\},
          \qquad r=\eta\ell,\qquad q=c_1\ell\Delta,\\
        M_i&=\mathcal G_{q,\Delta}
          \left(\sum_{j:\,d(i,j)\le\ell-R}H_j,\ H_i\right),
    \end{align*}
    where $\mathcal G_{q,\Delta}(H',W)$ is the normalized Gaussian
    filter of $W$ under $H'$.
    These actual truncated terms determine every projection of the
    corrected periodic hierarchy. If a unit ground vector $\psi$
    satisfies the source's local condition
    $\|(M_i-\langle\psi,M_i\psi\rangle)\psi\|\le J_b e^{-r}$,
    the hierarchy retains a unit vector $\phi$ at every node with
    $\|\phi-\psi\|\le2(3^{h+1}-1)e^{-r/8}$.
    Every finite orthonormal family retained by the final projection
    has at most $2\operatorname{tr}P_0$ members. Moreover,
    $\ell<\lceil R\rceil+A_M/\eta+1$.
    No localization estimate, pair-norm estimate, length inequality,
    or gap of an approximating Hamiltonian is an additional premise.
\end{theorem}
\begin{proof}\leanok
    \uses{thm:hastings_periodic_physical_metric,
        thm:hastings_periodic_physical_log_squared,
        thm:hastings_periodic_chosen_radius,
        thm:hastings05_metric_facld,
        thm:hastings_zero_range_exact}
    The physical ball count gives a uniform finite-range metric
    localization estimate for $R>0$. Thus the displayed $M_i$ are
    Hermitian, supported in their radius-$\ell$ balls, have norm at
    most $J$, and satisfy
    $\|\widetilde H_i-M_i\|\le JC e^{-r}\le J_b e^{-r}$.
    At $R=0$, all terms commute and both filters equal $H_i$ exactly.
    The maximum with $\lceil R\rceil$ ensures $\ell\ge R$.
    Since the block count is at most $M$, enlarging the chosen radius
    preserves its sufficient length inequality. Apply
    Theorem~\ref{thm:hastings_periodic_physical_log_squared} to these
    constructed terms and the source's local residual condition.
\end{proof}
